\documentclass[11pt,a4paper]{article}

% ── Packages ──────────────────────────────────────────────────────────────────
\usepackage[utf8]{inputenc}
\usepackage[T1]{fontenc}
\usepackage[french]{babel}
\usepackage{lmodern}
\usepackage{geometry}
\geometry{margin=2.5cm}
\usepackage{graphicx}
\usepackage{booktabs}
\usepackage{array}
\usepackage{multirow}
\usepackage{hyperref}
\usepackage{xcolor}
\usepackage{amsmath}
\usepackage{caption}
\usepackage{subcaption}
\usepackage{float}
\usepackage{parskip}

\hypersetup{colorlinks=true, linkcolor=blue, urlcolor=blue, citecolor=blue}

% ── Infos ─────────────────────────────────────────────────────────────────────
\title{
    \large ENSIMAG --- Projet Personnel\\[0.8em]
    \LARGE \textbf{Détection de Password Spraying}\\
    \large par Analyse Comportementale et Anomaly Detection
}
\author{Racim Zenati}
\date{Mars 2026}

% ══════════════════════════════════════════════════════════════════════════════
\begin{document}
\maketitle
\tableofcontents
\newpage

% ── 1. Introduction ───────────────────────────────────────────────────────────
\section{Introduction}

Le \textit{password spraying} est une attaque d'authentification consistant à
tester un petit nombre de mots de passe courants sur un grand nombre de comptes
distincts. Contrairement au brute force classique, cette approche vise à rester
sous les seuils de détection habituels (verrouillage de compte après $N$ échecs),
rendant son identification particulièrement difficile avec des méthodes classiques.

Ce projet propose un système de détection basé sur l'analyse comportementale des
flux d'authentification et les méthodes de détection d'anomalies non supervisées.
L'objectif est d'identifier des patterns suspects au niveau des adresses IP sources
sans nécessiter de données labellisées à l'entraînement.

\subsection*{Contributions}
\begin{itemize}
    \item Implémentation d'un pipeline de détection d'anomalies avec Isolation Forest
    \item Évaluation comparative sur deux datasets publics (CIC-IDS-2017, UNSW-NB15)
    \item Analyse des limites de la détection flux-par-flux face aux attaques discrètes
    \item Démonstration de l'apport de l'agrégation comportementale par IP
\end{itemize}

% ── 2. Données ────────────────────────────────────────────────────────────────
\section{Données}

\subsection{CIC-IDS-2017}

Le dataset CIC-IDS-2017 \cite{cicids2017} est produit par le Canadian Institute
for Cybersecurity. Il contient 2\,830\,743 flux réseau capturés sur une semaine,
représentant du trafic normal ainsi que 14 types d'attaques différentes.
Chaque flux est décrit par 78 features numériques extraites au niveau réseau
(durée, volume de paquets, flags TCP, etc.) et une colonne \texttt{Label}.

\begin{table}[H]
\centering
\caption{Distribution des labels dans CIC-IDS-2017}
\begin{tabular}{lrr}
\toprule
\textbf{Label} & \textbf{Occurrences} & \textbf{Proportion} \\
\midrule
BENIGN         & 2\,273\,097 & 80.3\% \\
DoS Hulk       & 231\,073    &  8.2\% \\
PortScan       & 158\,930    &  5.6\% \\
DDoS           & 128\,027    &  4.5\% \\
FTP-Patator    &   7\,938    &  0.3\% \\
SSH-Patator    &   5\,897    &  0.2\% \\
Autres         &  25\,781    &  0.9\% \\
\bottomrule
\end{tabular}
\end{table}

\subsection{UNSW-NB15}

Le dataset UNSW-NB15 \cite{unswnb15} est produit par l'université de New South
Wales. Il est fourni avec un split train/test prédéfini (175\,341 / 82\,332 lignes)
et contient 9 catégories d'attaques. Chaque flux est décrit par 34 features.
Ce dataset est particulièrement adapté à l'évaluation rigoureuse car le split
train/test est fixé indépendamment du ratio de contamination.

\begin{table}[H]
\centering
\caption{Distribution des catégories dans UNSW-NB15 (test set)}
\begin{tabular}{lrr}
\toprule
\textbf{Catégorie} & \textbf{Occurrences} & \textbf{Recall obtenu} \\
\midrule
Normal         & 37\,000 & ---    \\
Generic        & 18\,871 & 99.2\% \\
Exploits       & 11\,132 & 82.8\% \\
Fuzzers        &  6\,062 & 68.4\% \\
DoS            &  4\,089 & 91.8\% \\
Reconnaissance &  3\,496 & 67.2\% \\
Analysis       &    677  & 93.9\% \\
Backdoor       &    583  & 93.0\% \\
Shellcode      &    378  & 56.6\% \\
Worms          &     44  & 86.4\% \\
\bottomrule
\end{tabular}
\end{table}

% ── 3. Méthode ────────────────────────────────────────────────────────────────
\section{Méthode}

\subsection{Isolation Forest}

Isolation Forest \cite{iforest} est un algorithme de détection d'anomalies
non supervisé. Il construit un ensemble d'arbres de décision aléatoires en
isolant les observations : les anomalies, peu nombreuses et atypiques, sont
isolées en moins de découpages que les points normaux.

Le score d'anomalie d'un point $x$ est défini par :
\[
s(x, n) = 2^{-\frac{\mathbb{E}[h(x)]}{c(n)}}
\]
où $h(x)$ est la longueur de chemin moyenne dans les arbres et $c(n)$ est
la longueur moyenne attendue pour $n$ observations.

\subsection{Pipeline}

Le pipeline implémenté suit les étapes suivantes :

\begin{enumerate}
    \item \textbf{Chargement} des données brutes (CSV ou Parquet)
    \item \textbf{Nettoyage} : remplacement des valeurs infinies et NaN par 0
    \item \textbf{Sélection} des features numériques uniquement
    \item \textbf{Normalisation} par \texttt{StandardScaler} (moyenne 0, écart-type 1)
    \item \textbf{Entraînement} sur le trafic normal uniquement (non supervisé)
    \item \textbf{Scoring} : calcul du score d'anomalie sur l'ensemble du test set
    \item \textbf{Décision} : seuil basé sur le paramètre \texttt{contamination}
\end{enumerate}

\subsection{Agrégation comportementale par IP}

Pour améliorer la détection des attaques discrètes, nous proposons une
agrégation des flux par adresse IP source sur des fenêtres temporelles de 10
minutes. Les features calculées sont :

\begin{itemize}
    \item $n\_flows$ : nombre total de flux dans la fenêtre
    \item $n\_dst\_ports$ : nombre de ports de destination distincts
    \item $fail\_rate$ : taux d'échec des connexions
    \item $flows\_per\_min$ : débit de connexions par minute
    \item $port\_entropy$ : entropie des ports de destination
\end{itemize}

Cette approche s'inspire des méthodes Risk-Based Authentication (RBA)
\cite{rba} qui agrègent le comportement par utilisateur pour détecter
des connexions anormales.

% ── 4. Résultats ──────────────────────────────────────────────────────────────
\section{Résultats}

\subsection{Optimisation des hyperparamètres}

Une recherche sur grille a été conduite sur le dataset UNSW-NB15 pour les
paramètres \texttt{n\_estimators} $\in \{100, 200, 300\}$,
\texttt{max\_samples} $\in \{0.5, 0.8, 1.0\}$ et
\texttt{contamination} $\in \{0.3, 0.4, 0.499\}$.

\begin{table}[H]
\centering
\caption{Comparaison des modèles — meilleure configuration}
\begin{tabular}{lccc}
\toprule
\textbf{Modèle} & \textbf{F1} & \textbf{Recall} & \textbf{Precision} \\
\midrule
Isolation Forest & \textbf{0.743} & \textbf{82.8\%} & 100\% \\
One-Class SVM    & 0.736         & 85.6\%          & 100\% \\
Local Outlier Factor & 0.540    & 51.6\%          & 100\% \\
\bottomrule
\end{tabular}
\end{table}

Le meilleur modèle est \textbf{Isolation Forest} avec
\texttt{n\_estimators=100}, \texttt{max\_samples=1.0},
\texttt{contamination=0.499}, atteignant un F1 de 0.743 et un recall de
82.8\% sur le test set UNSW-NB15.

\begin{figure}[H]
    \centering
    \includegraphics[width=0.75\textwidth]{figures/fig2_comparaison_modeles.pdf}
    \caption{Comparaison F1 et Recall des trois modèles testés}
\end{figure}

\subsection{Distribution des scores}

\begin{figure}[H]
    \centering
    \includegraphics[width=0.75\textwidth]{figures/fig3_distribution_scores.pdf}
    \caption{Distribution des scores d'anomalie selon la classe réelle}
\end{figure}

La figure montre une bonne séparation entre trafic normal et attaques,
avec un chevauchement partiel expliquant les faux négatifs restants.

\subsection{Matrice de confusion}

\begin{figure}[H]
    \centering
    \includegraphics[width=0.5\textwidth]{figures/fig4_confusion_matrix.pdf}
    \caption{Matrice de confusion du modèle final sur UNSW-NB15}
\end{figure}

\subsection{Courbe Precision-Recall}

\begin{figure}[H]
    \centering
    \includegraphics[width=0.75\textwidth]{figures/fig5_precision_recall.pdf}
    \caption{Courbe Precision-Recall — modèle final}
\end{figure}

\subsection{Recall par catégorie d'attaque}

\begin{figure}[H]
    \centering
    \includegraphics[width=0.75\textwidth]{figures/fig6_recall_categories_unsw.pdf}
    \caption{Recall par catégorie d'attaque sur UNSW-NB15}
\end{figure}

\subsection{Analyse flux brut vs agrégation comportementale}

\begin{figure}[H]
    \centering
    \includegraphics[width=0.85\textwidth]{figures/fig1_recall_par_attaque.pdf}
    \caption{Recall par type d'attaque : CIC-IDS flux brut vs UNSW-NB15}
\end{figure}

L'analyse comparative sur CIC-IDS-2017 révèle une limite fondamentale de
l'approche flux-par-flux face aux attaques discrètes. Le tableau ci-dessous
résume les résultats :

\begin{table}[H]
\centering
\caption{Recall selon l'approche de détection sur CIC-IDS-2017}
\begin{tabular}{lcc}
\toprule
\textbf{Type d'attaque} & \textbf{Flux brut} & \textbf{Agrégation IP/10min} \\
\midrule
DoS GoldenEye  & 96.2\% & ---    \\
DoS Slowhttptest & 95.7\% & ---  \\
SSH-Patator    & 49.9\% & 50.0\% \\
FTP-Patator    &  1.3\% & \textbf{50.0\%} \\
PortScan       &  0.5\% & \textbf{91.1\%} \\
\bottomrule
\end{tabular}
\end{table}

L'agrégation comportementale par IP multiplie le recall de FTP-Patator
par \textbf{38$\times$} (de 1.3\% à 50\%), confirmant que les attaques
lentes et répétitives ne sont détectables qu'en analysant le comportement
global d'une IP sur une fenêtre temporelle.

% ── 5. Discussion ─────────────────────────────────────────────────────────────
\section{Discussion}

\subsection{Limites}

\textbf{Contamination fixe.} Le paramètre \texttt{contamination} contrôle
le seuil de décision directement. Sur CIC-IDS agrégé, fixer
\texttt{contamination} au taux réel d'attaques (31.6\%) produit un recall
artificiellement élevé. Les résultats UNSW-NB15, avec un split train/test
indépendant, sont plus fiables.

\textbf{Précision à 100\%.} La précision parfaite observée est un artefact
de la métrique sur ce type de dataset déséquilibré --- elle ne reflète pas
les performances en production où le taux de faux positifs est critique.

\textbf{Absence d'IP réelles dans CIC-IDS CSV.} Les fichiers CSV de
CIC-IDS ne contiennent pas les adresses IP sources, ce qui a nécessité
leur simulation pour l'agrégation comportementale.

\subsection{Perspectives}

\begin{itemize}
    \item Tester sur CIC-IDS-2018 qui contient des attaques brute force SSH/FTP plus récentes
    \item Intégrer une composante temporelle (LSTM ou sliding window) pour améliorer
          la détection des attaques lentes
    \item Comparer avec des méthodes supervisées (Random Forest, XGBoost) pour
          quantifier le coût du non-supervisé
\end{itemize}

% ── 6. Conclusion ─────────────────────────────────────────────────────────────
\section{Conclusion}

Ce projet démontre qu'Isolation Forest est une approche compétitive pour la
détection d'intrusions non supervisée, atteignant 82.8\% de recall et un F1
de 0.743 sur UNSW-NB15. L'analyse comparative sur CIC-IDS-2017 révèle que
l'agrégation comportementale par IP est indispensable pour détecter les
attaques discrètes de type password spraying : FTP-Patator passe de 1.3\%
à 50\% de recall avec cette approche.

Ces résultats motivent une architecture hybride combinant détection
flux-par-flux pour les attaques volumiques et agrégation comportementale
pour les attaques lentes, à l'image des systèmes RBA déployés en production.

% ── Références ────────────────────────────────────────────────────────────────
\begin{thebibliography}{9}

\bibitem{cicids2017}
Sharafaldin, I., Habibi Lashkari, A., \& Ghorbani, A. A. (2018).
\textit{Toward Generating a New Intrusion Detection Dataset and Intrusion Traffic Characterization}.
ICISSP 2018.

\bibitem{unswnb15}
Moustafa, N., \& Slay, J. (2015).
\textit{UNSW-NB15: A Comprehensive Data Set for Network Intrusion Detection Systems}.
MilCIS 2015.

\bibitem{iforest}
Liu, F. T., Ting, K. M., \& Zhou, Z. H. (2008).
\textit{Isolation Forest}.
ICDM 2008.

\bibitem{rba}
Freeman, J., \& Hayton, N. (2016).
\textit{Risk-Based Authentication: From Concept to Practice}.
USENIX Security 2016.

\end{thebibliography}

\end{document}
